\title{xLSTM: Extended Long Short-Term Memory}

\begin{abstract}
The foundational concepts of the Long Short-Term Memory (LSTM)—namely, the constant error carousel and the mechanism of gating—were established in the 1990s. LSTMs have since become robust mainstays in the deep learning landscape, underpinning multiple breakthroughs, and in particular, serving as the architecture behind the earliest Large Language Models (LLMs). The introduction of Transformer architectures, featuring parallelizable self-attention, represented a paradigm shift by eclipsing LSTM performance at scale. This paper investigates a fundamental question: What levels of language modeling performance can be achieved by scaling LSTMs to the order of billions of parameters, adopting recent advancements from the modern LLM landscape, and systematically addressing traditional LSTM constraints? Our approach introduces exponential gating mechanisms, complemented by appropriate normalization and stabilization methods. Furthermore, we re-envision the memory component of LSTMs, proposing: (i) sLSTM, which features a scalar-based memory and update mechanism coupled with revised memory mixing strategies; and (ii) mLSTM, which employs a fully parallelizable matrix-based memory with an associated covariance-driven update procedure. These LSTM enhancements are incorporated into residual block backbones, forming xLSTM blocks that are subsequently stacked to constitute the xLSTM architectures. The synergy of exponential gating and revamped memory designs enables xLSTM models to demonstrate compelling performance, rivaling the latest Transformers and State Space Models in both accuracy and scalability.\\
Code available at: \url{https://github.com/NX-AI/xlstm}
\end{abstract}

\section{Introduction}
\label{sec:intro}

Long Short-Term Memory (LSTM) concepts~\citep{Hochreiter:91,Hochreiter:97nips,Hochreiter:97}, such as gating mechanisms and the constant error carousel, were proposed as solutions to address the issue of vanishing gradients in recurrent neural networks~\citep{Hochreiter:91,Hochreiter:00book}:

\begin{align}
\label{eq:lstm_idea}
  c_t \ = \ f_t \ c_{t-1} \ + \ 
          i_t \  z_t \ , \quad  
          h_t \ = \ o_t \ \psi({c_t}) \ . 
\end{align}

The constant error carousel functions by adding updates to the cell state $c_{t-1}$ (shown in green) via cell inputs $z_t$, with this process being regulated by sigmoid gates (depicted in blue). Specifically, the input gate $i_t$ and forget gate $f_t$ govern the modifications to the cell state, while the output gate $o_t$ determines what portion of the memory cell is revealed as output, corresponding to the hidden state $h_t$. Before output gating, the cell state undergoes normalization or squashing through $\psi$, after which the gated value constitutes the hidden state.

LSTMs have achieved notable success across a range of fields \citep{Hochreiter:01,Hochreiter:07,Schmidhuber:15}, and held a dominant position in text generation until the emergence of Transformers in 2017~\citep{Vaswani:17}. Their utility has been established in numerous sequential data challenges, such as automatic text generation~\citep{Graves:13,Karpathy:15blog}, handwriting synthesis~\citep{Graves:13}, sequence-to-sequence translation tasks~\citep{Sutskever:14nips}, program evaluation~\citep{Zaremba:14arxiv}, generating image descriptions~\citep{Karpathy:15,Hossain:19}, source code synthesis~\citep{Karpathy:15blog}, modeling rainfall-runoff dynamics~\citep{Kratzert:18,Kratzert:19}, and constructing hydrological models for flood forecasting~\citep{Nearing:24}. In the context of reinforcement learning, LSTMs represent the most successful architectures for sequence processing—prominent examples include their deployment in the AlphaStar system for StarCraft II~\citep{Vinyals:17short}, OpenAI Five for Dota 2~\citep{OpenAI:19}, and magnetic controller models for nuclear fusion experiments~\citep{Degrave:22short}. A key strength of LSTMs lies in their capacity for abstraction, effectively identifying and storing semantic features within their memory components~\citep{Karpathy:15blog}, as shown by the discovery of specific units such as number and syntax neurons~\citep{Lakretz:19}, linguistic representation neurons~\citep{Bau:19}, and sentiment-sensitive neurons~\citep{Radford:17arxiv}. Despite evolving trends, LSTMs continue to be integral to significant contemporary applications~\citep{Degrave:22short,Nearing:24} and have demonstrated enduring value in machine learning research.

Despite their significant achievements, LSTMs face three primary drawbacks: (i) Inability to update prior storage choices. To illustrate, consider the 
{\it Nearest Neighbor Search} task (see Appendix~\ref{sec:appNearestNeighborSearch} for details): Given a reference vector, the sequence is processed to identify the vector most similar to the reference and return its corresponding value at the sequence's conclusion. The left side of Figure~\ref{fig:lstmProblems} displays the mean squared error for this task. The LSTM has difficulty replacing an already stored value upon encountering a more similar match later in the sequence. In contrast, our proposed xLSTM addresses this with exponential gating. (ii) Restricted memory capacity, whereby all information must be represented as scalar cell states. This shortcoming is demonstrated via the {\it Rare Token Prediction} task. The right side of Figure~\ref{fig:lstmProblems} presents token prediction perplexity on Wikitext-103~\citep{Merity:17} for groups categorized by token frequency. LSTM underperforms on low-frequency (rare) tokens because of its storage bottlenecks, whereas xLSTM alleviates this issue using matrix-based memory. (iii) Lack of parallel computation capability due to hidden-to-hidden connectivity, which necessitates strict sequential processing by mixing memory across time steps.

The constraints inherent to LSTM have driven the development of Transformers~\citep{Vaswani:17} for language modeling applications. This raises the question: by addressing these LSTM shortcomings and extending their capacity to match the scale of modern Large Language Models, what levels of performance could be realized in language modeling tasks?

\section{Extended Long Short-Term Memory}
\label{sec:xLSTM}

In addressing the limitations inherent to the standard LSTM, Extended Long Short-Term Memory (xLSTM) brings forth two principal innovations that build upon the foundational LSTM framework described in Equation~\eqref{eq:lstm_idea}. Specifically, these advancements — namely exponential gating and new memory organization schemes — result in the formulation of two novel variants within the LSTM paradigm: (i) the sLSTM (refer to Section~\ref{sec:sLSTM}), characterized by a single scalar memory, scalar update operations, and a mechanism for mixing memory states; and (ii) the mLSTM (see Section~\ref{sec:mLSTM}), which operates with matrix-based memory representations and utilizes a covariance (outer product) update strategy that is inherently compatible with parallel computation. Both the sLSTM and mLSTM models are augmented via exponential gating mechanisms. In pursuit of full parallelizability, mLSTM dispenses with memory mixing, thereby omitting recurrent connections among hidden states. Both architectures, sLSTM and mLSTM, are amenable to extensions involving multiple memory cells, where sLSTM, in particular, implements mixing across cells. Additionally, sLSTM may employ multiple heads; however, memory mixing is restricted to cells within an individual head, not across multiple heads. This architectural feature — introducing multiple heads in sLSTM in conjunction with exponential gating — represents a new form of memory mixing. In the case of mLSTM, the use of multiple heads does not differ conceptually from deploying multiple memory cells.

By embedding these novel LSTM variants within residual block modules, we form xLSTM blocks (refer to Section~\ref{sec:xLSTMblock}). Layering these xLSTM blocks in a residual fashion produces complete xLSTM architectures (refer to Section~\ref{sec:xLSTMarchitecure}). The overall structure and components of the xLSTM architecture are illustrated in Figure~\ref{fig:xlstm_sketch}.

\subsection{Review of the Long Short-Term Memory}
\label{sec:orgLSTM}

The initial LSTM framework~\citep{Hochreiter:91,Hochreiter:97nips,Hochreiter:97} proposed the scalar memory cell as its core unit for both information processing and retention, effectively mitigating the problem of vanishing gradients~\citep{Hochreiter:91,Hochreiter:00book} by utilizing what is termed the constant error carousel (i.e., the update of the cell state). Within the memory cell, three distinct gates exist: input gate, output gate, and forget gate, with the introduction of the forget gate attributed to \citet{Gers:00}. The temporal evolution of the LSTM cell is governed by the following update equations at each time step $t$:

\begin{align}
c_t \ &= \  f_t \ c_{t-1} \ + \ i_t \ z_t &  & &\text{cell state} \\
h_t  \ &= \ o_t \ \tilde{h}_t \ , & \tilde{h}_t \ &= \ \psi \left( c_t\right)
  &\text{hidden state} \\
z_t \ &= \ \varphi \left( \tilde{z}_t \right) \ , 
  &\tilde{z}_t \ &=  \ \mathbf{w}^\top_{z} \ \mathbf{x}_t \ + \
  r_{z}  h_{t-1} \ + \  b_{z} \ \
  &\text{cell input} \\
i_t \ &= \ \sigma \left( \tilde{i}_t  \right) \ , 
  &\tilde{i}_t \ &= \ \mathbf{w}^\top_{i} \ \mathbf{x}_t \ + \
  r_{i}  \ h_{t-1} \ + \  b_{i} \ \
  &\text{input gate} \\
f_t \ &= \ \sigma \left(  \tilde{f}_t \right) \ , 
  &\tilde{f}_t \ &= \ \mathbf{w}^\top_{f} \ \mathbf{x}_t  \ + \
  r_{f}  \ h_{t-1} \ + \  b_{f} \ \
  &\text{forget gate} \\
o_t \ &= \ \sigma \left( \tilde{o}_t \right) \ , 
  &\tilde{o}_t  \ &= \ \mathbf{w}^\top_{o} \ \mathbf{x}_t \ + \
  r_{o}  \ h_{t-1} \ + \  b_{o} \ \
  &\text{output gate} 
\end{align}

The above system of equations describes the internal computations of an LSTM unit. The cell state $c_t$ evolves by integrating the previous cell state $c_{t-1}$ modulated by the forget gate $f_t$ and adding the candidate memory $z_t$, which is regulated by the input gate $i_t$. The hidden state $h_t$ at time $t$ is determined as the element-wise product of the output gate $o_t$ and the transformed cell state $\tilde{h}_t$ through the activation function $\psi$. The intermediate cell input $z_t$ is derived by applying a nonlinearity $\varphi$ to its pre-activation $\tilde{z}_t$, which linearly combines the current input $\mathbf{x}_t$, the previous hidden state $h_{t-1}$ (weighted by $r_z$), and a bias term $b_z$. Similarly, each gate $i_t$, $f_t$, and $o_t$ is computed via a sigmoid nonlinearity $\sigma$ applied to their respective affine transformations $\tilde{i}_t$, $\tilde{f}_t$, and $\tilde{o}_t$, each constructed in an analogous manner from the current input, prior hidden state, and respective biases and weights. This structure facilitates the selective integration and transmission of information over sequential data.

The input weight vectors $\mathbf{w}_{z}$, $\mathbf{w}_{i}$,
$\mathbf{w}_{f}$, and $\mathbf{w}_{o}$ are associated with the connections from the inputs $\mathbf{x}_t$ to the cell input, input gate, forget gate, and output gate, respectively. The recurrent weights $r_{z}$, $r_{i}$, $r_{f}$, and $r_{o}$ link the previous hidden state $h_{t-1}$ to the cell input, input gate, forget gate, and output gate, correspondingly. The parameters $b_{z}$, $b_{i}$, $b_{f}$, and $b_{o}$ represent the respective bias terms.
The activation functions for the cell input and hidden state are denoted by $\varphi$ and $\psi$ (commonly implemented as $\tanh$).
$\psi$ plays the role of scaling or constraining the cell state, which could otherwise grow without bound. The activation function for all gates is the sigmoid function, given by $\sigma \left( x \right) = 1/(1+\exp(-x))$.
In subsequent models, several scalar memory cells $\mathbf{c}_t \in \mathbb{R}^{d}$ were aggregated into a vector, enabling recurrent weight matrices $\mathbf{R} \in \mathbb{R}^{d \times d}$ to combine outputs across the different memory cells~\citep{Greff:15} (refer to Appendix~\ref{sec:apporgLSTM} for further discussion).
Findings from ablation studies indicated that each part of the memory cell contributes essentially to its operation~\citep{Greff:15}.

\subsection{sLSTM}
\label{sec:sLSTM}

To enable LSTMs to reconsider their storage choices, we incorporate exponential gates (depicted in red), along with mechanisms for normalization and stabilization. Specifically, the input and forget gates may utilize exponential activation functions. For the normalization process, a normalizer state is established to accumulate the sum of the input gate multiplied by each subsequent forget gate.

The scalar sLSTM forward pass is:

\begin{align}
\label{eq:slstmforward}
c_t \ &= \  f_t \ c_{t-1} \ + \ i_t \ z_t & & 
 &\text{cell state} \\
n_t \ &= \  f_t \ n_{t-1} \ + \ i_t  & & 
 &\text{normalizer state} \\
h_t \ &= \ o_t \ \tilde{h}_t \ , 
  &\tilde{h}_t \ &= \ c_t / n_t
&\text{hidden state} \\
z_t \ &= \ \varphi \left( \tilde{z}_t \right) \ , 
  &\tilde{z}_t \ &=  \ \mathbf{w}^\top_{z} \ \mathbf{x}_t \ + \
  r_{z}  h_{t-1} \ + \  b_{z}
  &\text{cell input} \\
i_t \ &= \ \!\exp \left( \tilde{i}_t  \right) \ , 
  &\tilde{i}_t \ &= \ \mathbf{w}^\top_{i} \ \mathbf{x}_t \ + \
  r_{i}  \ h_{t-1} \ + \  b_{i}
  &\text{input gate} \\
f_t \ &= \ \sigma \left(  \tilde{f}_t \right) \ \text{OR} \
\!\exp \left(  \tilde{f}_t \right) \ ,
  &\tilde{f}_t \ &= \ \mathbf{w}^\top_{f} \ \mathbf{x}_t  \ + \
  r_{f}  \ h_{t-1} \ + \  b_{f}
  &\text{forget gate} \\
o_t \ &= \ \sigma \left( \tilde{o}_t \right) \ , 
  &\tilde{o}_t  \ &= \ \mathbf{w}^\top_{o} \ \mathbf{x}_t \ + \
  r_{o}  \ h_{t-1} \ + \  b_{o}
  &\text{output gate} 
\end{align}

We adapt the conventional LSTM gating mechanisms—which incorporate dependencies on inputs and/or hidden states along with a bias term—into the proposed architectures. Given that exponential activation functions are prone to producing extremely large outputs, risking numerical overflows, we introduce an auxiliary state, \mbox{$m_t$~\citep{Milakov:18arxiv}}, to enhance the stability of the gates:

\begin{align}
\label{eq:slstmstabil}
m_t \ &= \ \max \left( \log ( f_t ) + m_{t-1},\ \log ( i_t ) \right) &\text{stabilizer state} \\
i'_t \ &= \ \!\exp \big( \log ( i_t ) - m_t \big) \ = \ \!\exp \left( \tilde{i}_t - m_t \right) \ \, 
  &\text{stabil. input gate} \\
f'_t \ &= \ \!\exp \big( \log ( f_t ) + m_{t-1} - m_t \big) \ 
  &\text{stabil. forget gate}
\end{align}

As demonstrated in Appendix~\ref{sec:appsLSTM}, substituting $f_t$ with $f'_t$ and $i_t$ with $i'_t$ during the forward computation does not alter the overall network output or affect the gradients of the loss with respect to its parameters.

\paragraph{New Memory Mixing.} Similar to the original LSTM architecture (see Appendix~\ref{sec:appsLSTM}), the sLSTM is capable of utilizing several memory cells. The existence of multiple cells allows for memory mixing through the recurrent links $\mathbf{R}_{\mathbf{z}}$, $\mathbf{R}_{\mathbf{i}}$, $\mathbf{R}_{\mathbf{f}}$, and $\mathbf{R}_{\mathbf{o}}$, which connect the hidden state $\mathbf{h}$ to the cell input $\mathbf{z}$ and the various gates $\mathbf{i}$, $\mathbf{f}$, and $\mathbf{o}$. In this context, memory mixing acquires a novel dimension owing to the adoption of exponential gating mechanisms. Furthermore, the latest formulation of sLSTM permits multiple heads, facilitating memory mixing within the boundaries of each head, but not between different heads. Consequently, the integration of sLSTM heads and exponential gating introduces a fundamentally new approach to memory mixing.

\subsection{mLSTM}
\label{sec:mLSTM}

In order to boost the memory capacity of LSTMs, we expand the standard scalar memory cell $c \in \mathbb{R}$ to a full matrix $\mathbf{C} \in \mathbb{R}^{d \times d}$. This modification allows for memory access through matrix multiplication operations. Specifically, at each time step $t$, the architecture is designed to store a pair of $d$-dimensional vectors: the key $\mathbf{k}_t \in \mathbb{R}^d$ and the value $\mathbf{v}_t \in \mathbb{R}^d$ (adopting the nomenclature from Transformer models). Subsequently, at a later timestep $t + \tau$, it is desired that the value $\mathbf{v}_t$ can be accessed or retrieved using an associated query vector $\mathbf{q}_{t+\tau} \in \mathbb{R}^d$. This setup mirrors the principles found in Bidirectional Associative Memories (BAMs) as introduced in \citep{Kohonen:72,Anderson:72,Nakano:72,Anderson:77}.

The covariance update rule~\citep{Sejnowski:77,Dayan:91} for storing a key-value pair is

\begin{align}
\mathbf{C}_t \ &= \ \mathbf{C}_{t-1} \ + \ \mathbf{v}_{t} \ \mathbf{k}_t^\top \ .
\end{align}

We consider the scenario where a layer normalization step is applied prior to projecting the input data into key and value representations, resulting in these vectors having a mean of zero. The rule for updating the covariance is theoretically optimal~\citep{Dayan:91} for maximizing the distinguishability among retrieved binary codes, which aligns with achieving the highest possible signal-to-noise ratio. This distinctiveness can be further increased, though at the cost of increased (quadratic) computational resources, when restricting retrieval mechanisms to pairwise dependencies—akin to attention mechanisms~\citep{Krotov:16,Krotov:17,Ramsauer:21}. Moreover, the covariance update algorithm corresponds to the mechanics of Fast Weight Programmers~\citep{Schmidhuber:92ncfastweights,Schlag:21}, which have subsequently been enhanced by introducing a constant decay factor applied to $\mathbf{C}_{t-1}$ and a constant learning factor for the update $\mathbf{v}_{t} \mathbf{k}_t ^\top$~\citep{Ba:16}. Consistent with this approach, we embed the covariance-based update procedure within the LSTM model architecture, where the forget gate realizes the decay mechanism, the input gate modulates the learning rate, and the output gate adjusts the amplitude of the retrieved result.

In this matrix memory framework, the normalizer state is computed as a sum of key vectors, with each key vector multiplied by both the input gate and the cumulative product of all subsequent forget gates. The normalizer thus maintains information regarding the overall gate activation magnitudes. To prevent situations where the dot product of the query and normalizer state approaches zero, we take the absolute value of this dot product and apply a lower bound threshold, commonly set to 1.0, as suggested in prior work~\citep{Sun:23arxiv}. The forward computation in mLSTM is given as follows:

\begin{align}
\label{eq:mlstm_recurrent_begin}
\mathbf{C}_t \ &= \  f_t \ \mathbf{C}_{t-1} \ + \ 
  i_t \ \mathbf{v}_t \ \mathbf{k}_t^\top & &  &\text{cell state} \\
\mathbf{n}_t \ &= \  f_t \ \mathbf{n}_{t-1} \ + \ 
  i_t \ \mathbf{k}_t & &  &\text{normalizer state} \\
\mathbf{h}_t  \ &= \ \mathbf{o}_t \ \odot \ \tilde{\mathbf{h}}_t
\ , \qquad \qquad 
\tilde{\mathbf{h}}_t \ = \   \mathbf{C}_t \mathbf{q}_t \ / \ 
  \max \left\{ |\mathbf{n}_t^\top \mathbf{q}_t|, 1 \right\} 
   &&&\text{hidden state} \\
\mathbf{q}_t \ &= \ \mathbf{W}_q \ \mathbf{x}_t \ + \ \mathbf{b}_q  & &  &\text{query input} \\
\mathbf{k}_t \ &= \ \frac{1}{\sqrt{d}} \mathbf{W}_k \ \mathbf{x}_t \ + \ \mathbf{b}_k  & &  &\text{key input} \\
\mathbf{v}_t \ &= \ \mathbf{W}_v \ \mathbf{x}_t \ + \ \mathbf{b}_v  & &  &\text{value input} \\
i_t \ &= \ \!\exp \!  \! \left( \tilde{i}_t  \right) \ , 
 \qquad \qquad \ \ \ \ \,
  \tilde{i}_t \ = \ \mathbf{w}^\top_{i} \ \mathbf{x}_t \ + \  b_{i} \ \ \ 
  &&&\text{input gate} \\
f_t \ &= \ \sigma \!  \left(  \tilde{f}_t \right) \ \text{OR} \ \!\exp \!  \! \left(  \tilde{f}_t \right) , \ \ \: \!
\tilde{f}_t \ = \ \mathbf{w}^\top_{f} \ \mathbf{x}_t  \ + \
  b_{f}
  &&& \text{forget gate} \\
\label{eq:mlstm_recurrent_end}
\mathbf{o}_t \ &= \ \sigma \left( \tilde{\mathbf{o}}_t \right) \ , \qquad \qquad \quad \ \ \ \,
  \tilde{\mathbf{o}}_t  \ = \ \mathbf{W}_{\mathbf{o}} \ \mathbf{x}_t \ + \
  \mathbf{b}_{\mathbf{o}}
  &&&\text{output gate} 
\end{align}

The update equation for the cell state, $\mathbf{C}_t$, is a sum of the previous cell state scaled by the forget gate $f_t$ and the outer product between the input-gated value $\mathbf{v}_t$ and the key $\mathbf{k}_t$, where both are modulated by $i_t$. The normalizer state $\mathbf{n}_t$ evolves in a similar fashion, updating as a combination of its prior value weighted by $f_t$ and the addition of $i_t \mathbf{k}_t$. The hidden state $\mathbf{h}_t$ is calculated as the elementwise product between the output gate $\mathbf{o}_t$ and a normalized version $\tilde{\mathbf{h}}_t$; the latter normalizes $\mathbf{C}_t \mathbf{q}_t$ by the larger of either one or the absolute value $|\mathbf{n}_t^\top \mathbf{q}_t|$. At each step, the input $\mathbf{x}_t$ is projected to produce the query $\mathbf{q}_t$, key $\mathbf{k}_t$, and value $\mathbf{v}_t$ vectors via linear transformations plus biases. The input gate $i_t$ is determined by an exponentiated linear combination of $\mathbf{x}_t$, while the forget gate $f_t$ is a nonlinear function (either sigmoid or exponential) applied to its own affine transformation of $\mathbf{x}_t$. Finally, the output gate $\mathbf{o}_t$ uses a sigmoid applied to a corresponding affine function of $\mathbf{x}_t

mLSTM, similar to traditional LSTM, is capable of supporting several memory cells. In the case of mLSTM, employing multiple heads is functionally identical to using multiple cells because memory states do not interact with each other. To ensure stability in the exponential gating mechanisms of mLSTM, we apply the stabilization methods developed for sLSTM (refer to Equation~\ref{eq:slstmstabil}). Given the lack of memory mixing in mLSTM, its recurrence relation permits a reformulation that allows for parallel computation. Further elaboration can be found in Appendix~\ref{sec:appmLSTM}.

\subsection{xLSTM Architecture}
\label{sec:xLSTMarch}

\paragraph{xLSTM Blocks.}
\label{sec:xLSTMblock}
An xLSTM block aims to encode previous information in a non-linear manner within a high-dimensional feature space, facilitating the discrimination between diverse historical states or contexts. Effective separation of these prior states is essential for accurately forecasting subsequent elements in a sequence, such as the upcoming token. Our approach is motivated by Cover’s Theorem~\citep{Cover:65}, which posits that patterns embedded non-linearly into a space of higher dimensionality are more likely to achieve linear separability compared to their original embedding. We investigate two types of residual block configurations: (i) The residual block employing a post up-projection mechanism (similar to those found in Transformers), which first non-linearly aggregates past information in its initial space, then executes a linear mapping to a higher-dimensional representation, follows with a non-linear activation, and finally projects the features back to the original space via another linear mapping. This configuration is illustrated in the left section of Figure~\ref{fig:Backbones} and the third column of Figure~\ref{fig:xlstm_sketch}. A more granular schematic is provided in Figure~\ref{fig:appsLSTM_detailed} located in the appendix. (ii) A residual block based on pre up-projection principles (akin to State Space Models), which performs a linear transformation into a high-dimensional domain,

The model first performs a nonlinear compression of previous information into a high-dimensional representation, which is subsequently projected linearly back to the input space. In cases where an xLSTM block includes an sLSTM module, we employ the up-projection block after this step. Conversely, if the xLSTM incorporates an mLSTM module, the up-projection block is introduced beforehand, reflecting the increased memory capacity available in the high-dimensional context. For further clarification, consult the left side of Figure~\ref{fig:Backbones}, the third column in Figure~\ref{fig:xlstm_sketch}, or refer to Figure~\ref{fig:appsLSTM_detailed} provided in the appendix.

\paragraph{xLSTM Architecture. }
\label{sec:xLSTMarchitecure}
The xLSTM model is formed by stacking fundamental modules in a residual manner~\citep{Srivastava:15,He:16}. Our approach utilizes the pre-LayerNorm~\citep{Ba:16b} residual architecture, which is a standard choice in current Large Language Models. Refer to the final column in Figure~\ref{fig:xlstm_sketch}.

\subsection{Memory and Speed Considerations}

In contrast to Transformers, xLSTM networks exhibit linear computational requirements and maintain fixed memory usage regardless of the input sequence length. Owing to the compressive nature of xLSTM memory, these networks are particularly advantageous for deployment in industrial environments and on edge devices.

The memory mechanism in mLSTM is parameter-free, yet its reliance on a $d \times d$ matrix structure and corresponding $d \times d$ update operations incurs high computational cost. This design choice reflects a balance between the system’s memory capacity and computational demands. However, because these calculations can be parallelized effectively on GPUs, the impact on actual runtime is minimal.

Although mLSTM can be parallelized similarly to FlashAttention~\citep{Dao:22,Dao:23} or GLA~\citep{Yang:23arxiv}, sLSTM lacks parallelizability because of its memory mixing stemming from hidden-to-hidden connections. Nonetheless, we engineered an efficient CUDA version that employs GPU memory optimizations down to register usage, resulting in sLSTM running at no more than twice the computation time of mLSTM.

\section{Related Work}
\label{sec:related}

\paragraph{Linear Attention.} Various strategies have been introduced to address the quadratic dependence on context length inherent in the Transformer’s attention mechanism and to achieve linear scalability instead. The Synthesizer generates learned, synthetic attention weights without utilizing direct token-to-token interactions~\citep{Tay:20arxiv}. Linformer achieves self-attention by employing a low-rank matrix to provide a linear approximation~\citep{Wang:20arxiv}. Linear Transformer presents a linearized version of the attention computation~\citep{Katharopoulos:20}. Performer utilizes a method based on positive orthogonal random features to obtain a linear approximation of the attention softmax~\citep{Choromanski:21}. In other works, such as Structured Global Convolution (SGConv)~\citep{Li:22} and the Hyena Hierarchy~\citep{Poli:23}, the traditional attention mechanism is replaced with efficient long convolution operations.

\paragraph{State Space Models.} In recent years, State Space Models (SSMs) have gained considerable traction owing to their linear scalability with respect to context length and competitive results when compared to Transformers. The Structured State Space sequence model (S4)~\citep{Gu:21} marked one of the earliest contributions in this area, which was subsequently extended by models such as the Diagonal State Space (DSS)~\citep{Gupta:22}, Gated State Space (GSS)~\citep{Mehta:22}, S5~\citep{Smith:22}, Bidirectional Gated SSM (BiGS)~\citep{Wang:22}, H3~\citep{Fu:23}, and Mamba~\citep{Gu:24arxiv}.

\paragraph{Recurrent Neural Networks.} Recurrent Neural Networks (RNNs) have emerged as an alternative to Transformers and attention mechanisms, primarily due to their ability to process context with linear complexity. Deep Linear Recurrent Units (LRUs) incorporated into RNNs have delivered competitive outcomes in language modeling tasks~\citep{Orvieto:23, De:24arxiv}, similar to the promising results obtained by Hierarchically Gated Linear RNNs (HGRN)~\citep{Qin:23} and its successor HGRN2~\citep{Qin:24arxiv}. Among the most prominent RNN-based approaches for large-scale language models is RWKV~\citep{Peng:23arxivshort,Peng:24arxivshort}, which rivals Transformers in performance.

\paragraph{Gating.}
The mechanism of gating, a foundational concept introduced by LSTM, has been revisited and adopted in various recent methodologies. Notably, gating plays a central role in models such as HGRN~\citep{Qin:23}, HGRN2~\citep{Qin:24arxiv}, Gated Linear Attention (GLA)~\citep{Yang:23arxiv}, Gated State Space (GSS) models~\citep{Mehta:22}, Bidirectional Gated SSM (BiGS)~\citep{Wang:22}, Moving Average Equipped Gated Attention (MEGA)~\citep{Ma:22}, as well as RWKV~\citep{Peng:23arxivshort} and Mamba~\citep{Gu:24arxiv}.

\paragraph{Covariance Update Rule.} In order to increase the storage capacity, we integrated a matrix memory governed by a covariance update rule within the mLSTM cell. Several other approaches also utilize a similar update mechanism, including Fast Weight Programmers \citep{Schmidhuber:92ncfastweights,Schlag:21}, RWKV-5 and RWKV-6~\citep{Peng:24arxivshort}, Retention~\citep{Sun:23arxiv}, Linear Transformer~\citep{Katharopoulos:20}, and HGRN2~\citep{Qin:24arxiv}.

\paragraph{Most Related.} The models most akin to xLSTM from a conceptual standpoint are Retention~\citep{Sun:23arxiv}, RWKV~\citep{Peng:23arxivshort,Peng:24arxivshort}, and HGRN2~\citep{Qin:24arxiv}. These approaches incorporate principles such as matrix memory and/or the use of gating mechanisms. Nonetheless, unlike the proposed sLSTM, these methods lack the capacity for memory mixing. The inclusion of memory mixing is critical for effectively addressing state tracking tasks, positioning LSTMs as more expressive than both State Space Models (SSMs) and Transformers~\citep{Merrill:24,Deletang:23}. Tasks such as code evaluation and the tracking of entities within extended narratives necessitate robust state tracking capabilities.

\paragraph{Residually Stacking Architectures.} Much like most modern large-scale deep neural networks, xLSTM architectures employ a residual stacking approach in their design, utilizing foundational components~\citep{Srivastava:15,He:16}.

This residual stacking technique has been fundamental to the advancement of deep convolutional networks~\citep{He:16} as well as the development of Transformer models~\citep{Vaswani:17}. The Transformer framework, in particular, serves as the core technology driving prominent Large Language Models (LLMs) such as GPT-3~\citep{Brown:20short}, ChatGPT~\citep{Schulman:22short}, GPT-4~\citep{Achiam:23gpt4short}, Megatron-LM~\citep{Shoeybi:19}, Gopher~\citep{Rae:21short}, ERNIE 3.0 Titan~\citep{Wang:21arxivshort}, GLaM~\citep{Du:21glamshort}, Chinese M6~\citep{Lin:21}, multilingual AlexaTM 20B~\citep{Soltan:22}, OPT~\citep{Zhang:22opt}, Chinchilla~\citep{Hoffmann:22short}, BLOOM~\citep{Scao:22arxivshort}, GLM-130B~\citep{Zeng:22arxivshort}, LaMDA~\citep{Thoppilan:22short}, PaLM~\citep{Chowdhery:22short}, Llama~\citep{Touvron:23arxiv}, and Gemini~\citep{GeminiTeam:23arxiv,Reid:24arxiv}.

\section{Experiments}
\label{sec:Exp}

In this section, we carry out a series of experiments to assess xLSTM and benchmark it against established approaches, concentrating primarily on the domain of language modeling. Section~\ref{sec:ExpSyntethic} explores the unique properties of xLSTM by subjecting it to targeted synthetic benchmarks. 
Subsequently, in Section~\ref{sec:ExpComparison}, we analyze the validation perplexity scores of a range of contemporary language models, all trained on 15B tokens derived from SlimPajama~\citep{Soboleva:23}. Here, we additionally conduct ablation studies focused on xLSTM to disentangle its various components. We further examine how different approaches scale, following the methodology outlined in \citet{Kaplan:20} and \citet{Brown:20short}.

Section~\ref{sec:ExpLanguage} is dedicated to a more comprehensive language modeling study, comparing xLSTM and the highest performing models from Section~\ref{sec:ExpComparison} after training them with an extended corpus of 300B tokens from SlimPajama~\citep{Soboleva:23}. In this part, we begin by evaluating the models’ abilities to generalize to longer contexts. We then assess model effectiveness using validation perplexity measures and a suite of downstream tasks~\citep{Sutawika:24}. Furthermore, we test the methods across 571 individual text domains from the PALOMA benchmark dataset~\citep{Magnusson:23arxivshort}. Lastly, we revisit the topic of model scaling, this time examining trends as the training dataset is expanded by a factor of twenty.

Across all our experiments, we adopt the notation xLSTM[$a$:$b$] to indicate that the composition of xLSTM blocks has an $a/b$ split between mLSTM-based and sLSTM-based blocks, respectively. For instance, xLSTM[7:1] refers to a configuration where, in a group of eight blocks, seven are built using mLSTM while only one utilizes sLSTM. Applying this to a typical total of 48 blocks results in 42 mLSTM-based blocks and 6 sLSTM-based ones. Additionally, all experimental setups employ up-projection blocks: for mLSTM blocks these are placed before (pre), and for sLSTM blocks these are placed after (post) the primary layers.

\subsection{Synthetic Tasks and Long Range Arena}
\label{sec:ExpSyntethic}
We begin by evaluating the impact of xLSTM’s exponential gating combined with memory mixing on tasks involving formal languages~\citep{Deletang:23}. Subsequently, we investigate the capabilities of xLSTM’s matrix memory through the Multi-Query Associative Recall benchmark~\citep{Arora:23arxiv}. Finally, we measure xLSTM’s capacity to handle extended input sequences by benchmarking its performance on the Long Range Arena suite~\citep{Tay:21}.

\paragraph{Evaluation of xLSTM's Exponential Gating Incorporating Memory Mixing.} We investigate the capabilities of xLSTM's novel exponential gating mechanism with integrated memory mixing, which is designed to facilitate state tracking tasks~\citep{Merrill:24, Merrill:23}. Our approach extends the set of formal language benchmarks from \citet{Deletang:23}, modifying them to support training across variable input lengths and to assess extrapolation performance. A comprehensive overview of each task, as well as additional empirical outcomes, is available in Appendix~\ref{sec:appExpSynthetic-formalLang}. 
The performance of xLSTM is assessed against a variety of baseline architectures, including Transformers, State Space Models, and conventional Recurrent Neural Networks. We report accuracy exclusively on tokens integral to the task objectives, standardizing the metric on a scale from 0 (equivalent to chance) to 1 (optimal performance). Across these benchmarks, we systematically compare two-block configurations: xLSTM[0:1] (i.e., the pure sLSTM variant), xLSTM[1:0] (i.e., employing only mLSTM), xLSTM[1:1], along with Llama, Mamba, RWKV, Retention, Hyena, LSTM, and an LSTM integrated within Transformer blocks (LSTM (Block)). Figure~\ref{fig:formal-main} illustrates the comparative findings. Notably, models lacking memory mixing—such as standard Transformers or State Space Models without state tracking mechanisms—are unable to resolve challenges involving regular grammars, such as the parity task. This observation reinforces existing claims that, in terms of expressiveness, Transformers and State Space Models fall short compared to RNNs~\citep{Merrill:24, Merrill:23, Deletang:23}.

\paragraph{Test of xLSTM's Memory Capacities on Associative Recall Tasks.} This study investigates the memory abilities of xLSTM’s novel matrix memory on the Multi-Query Associative Recall task~\citep{Arora:23arxiv}. In this setting, sequences are constructed by randomly sampling key-value pairs from an extensive vocabulary; these pairs need to be stored and subsequently recalled. To increase the challenge relative to the standard formulation, we raise the number of key-value associations up to 256 and extend the possible context length to 2048, thereby enabling a comprehensive evaluation of various models’ memory capacities. Our comparison focuses on 2-block architectures from Llama, Mamba, RWKV-5, RWKV-6, xLSTM[1:1], and xLSTM[1:0]. The metric used is the recall accuracy of key-value pairs. Transformers such as Llama, possessing exponentially large memory in relation to their coding dimension~\citep{Ramsauer:21}, serve as a benchmark for this task. Performance outcomes are summarized in Figure~\ref{fig:mqar-main}. Among models not based on the Transformer architecture, xLSTM[1:1] achieves the strongest performance, including for smaller model variants. Notably, the inclusion of the sLSTM block does not hinder memory but instead augments it, especially visible when handling 256 key-value pairs, the most demanding scenario. Further analysis can be found in Appendix~\ref{sec:appExpSynthetic-mqar}, with extrapolation experiments revealing that the xLSTM’s strengthened memory design supports generalization to contexts longer than those present during training.

\paragraph{Test of xLSTM's Long Context Capabilities on Long Range Arena.} To evaluate how well xLSTM manages extended sequences and broad contextual information, we benchmark it against alternative approaches using the Long Range Arena~\citep{Tay:21}. Across all evaluated tasks, xLSTM consistently achieves superior results, indicating that its architecture is highly effective and robust for addressing various challenges associated with long-context sequence modeling.

For more details, see Appendix~\ref{sec:appExpSynthetic-lra}.

\subsection{Method Comparison and Ablation Study}
\label{sec:ExpComparison}

In order to investigate the central question posed in this study—namely, how effective the proposed LSTM variants are when applied to large-scale language modeling—we conduct experiments using xLSTMs, Transformers, State Space Models, and additional approaches. All models are trained auto-regressively on 15B tokens sourced from SlimPajama. We evaluate the resulting models on the validation dataset and carry out ablation analyses specifically on the xLSTMs.

\paragraph{Comparing xLSTM to Other Methods.} To provide a fair comparison, we trained all models using 15B tokens drawn from the SlimPajama dataset~\citep{Soboleva:23}. The evaluation criterion for these models is perplexity measured on the designated validation set. The suite of models assessed includes: the xLSTM (our proposed approach), GPT-3 (Transformer)~\citep{Brown:20short}, Llama (Transformer)~\citep{Touvron:23arxiv}, H3 (SSM)~\citep{Fu:23}, Mamba (SSM)~\citep{Gu:24arxiv}, RWKV-4, RWKV-5, RWKV-6 (RNN-based)~\citep{Peng:23arxivshort, Peng:24arxivshort}, GLA (linear Transformer)~\citep{Yang:23arxiv}, HGRN and HGRN2 (RNNs)~\citep{Qin:23, Qin:24arxiv}, RetNet (linear Transformer)~\citep{Sun:23arxiv}, Hyena (linear Transformer)~\citep{Poli:23}, as well as variations xLSTM[1:0] and xLSTM[7:1] (notations detailed in Section~\ref{sec:xlstm_blocks_notation}). Training was conducted using mixed-precision techniques; however, for models such as RWKV-5, RWKV-6, GLA, and HGRN2, PyTorch’s automated mixed-precision capabilities were not applied (additional information is available in Appendix Section~\ref{sec:appGenTrainProc}). 

We group the considered methods into three main categories: (a) Transformers, (b) State Space Models (SSMs), and (c) Recurrent Neural Networks (RNNs), including linear Transformers, which replace traditional Transformer attention with linear operations. All architectures are configured to match a 350M-parameter GPT-3 baseline, consisting of a 1024-dimensional embedding space and 24 residual layers. Notably, GPT-3 uniquely ties token and output embedding weights, resulting in a reduced overall parameter count relative to others. According to the empirical results summarized in Table~\ref{tab:spaj15b_model_results}, xLSTM achieves superior validation perplexity over all competing baselines. Additional experimental details are provided in Appendix~\ref{sec:appExpComparison}. Scaling properties for this evaluation, as depicted in Figure~\ref{fig:temp_sclaw15B}, further suggest that xLSTM's advantage persists or widens with larger parameter regimes.

\paragraph{Ablation Studies.}
As shown in Table~\ref{tab:spaj15b_model_results} and Figure~\ref{fig:temp_sclaw15B}, xLSTM demonstrates strong performance in language modeling when trained with 15B tokens from the SlimPajama dataset. To isolate the impact of various modifications from the original LSTM to the xLSTM, we sequentially adapt a standard LSTM architecture into xLSTM. The process begins by embedding LSTM layers within residual backbones that use pre-LayerNorm. The next stage involves incorporating a post up-projection block. In the final step, we introduce exponential gating mechanisms along with matrix memory to complete the transformation.

The results are shown in Table~\ref{tab:ablstudies} (top). 

Ablation studies indicate that the significant gains in performance stem from the combined effects of exponential gating and the matrix memory. Given the pivotal role of gating in both RNNs and State Space Models, we further investigate the impact of various gating strategies. As shown in Table~\ref{tab:ablstudies} (bottom), our results demonstrate that making each gate trainable and dependent on the input leads to consistent incremental improvements. Further details regarding the ablations of specific backbone components can be found in Appendix~\ref{sec:appAblDetails}.

\subsection{xLSTM as Large Language Model}
\label{sec:ExpLanguage}

This research concludes with comprehensive language modeling experiments aimed at evaluating xLSTM’s capabilities as a large language model (LLM). To this end, we significantly expand the training dataset, employing 300B tokens sourced from SlimPajama. For consistency, this token count matches that utilized in other works such as Mamba~\citep{Gu:24arxiv} and Griffin~\citep{De:24arxiv}.

Our analysis benchmarks xLSTM against RWKV-4, Llama, and Mamba, each exemplifying a unique methodological approach as detailed in Section~\ref{sec:ExpComparison}. RWKV-4 serves as the RNN baseline due to suitable training precision for RWKV-5, RWKV-6, and HGRN2 only being established after commencement of the 300B token runs (see Appendix~\ref{sec:appGenTrainProc}). We explore multiple parameter scales (125M, 350M, 760M, and 1.3B), examining not only the models’ capacity for sequence length extrapolation but also their accuracy on the validation data. Further, we evaluate these models on various downstream applications, quantify their language modeling performance across the 571 textual categories of the PALOMA benchmark, and systematically study their adherence to empirical scaling laws.

\paragraph{Sequence Length Extrapolation.}
\label{sec:spaj300B_ctx_extrapolation}
To begin, we evaluate how well large models with 1.3B parameters—specifically xLSTM, RWKV-4, Llama, and Mamba—generalize to longer sequence lengths. Each model is trained with a context length of 2048, and then assessed on increasingly larger contexts, up to 16384 tokens. The outcomes are depicted in Figure~\ref{fig:spaj300B_seq_extrapolation}. Notably, xLSTM consistently achieves lower perplexity when handling extended contexts, outperforming the alternative approaches in this regard.

\paragraph{Validation Perplexity and Downstream Tasks.}
\label{sec:spaj_downstream_eval}
In addition, across all model scales, we assess the effectiveness of xLSTM, RWKV-4, Llama, and Mamba architectures by measuring their performance on the SlimPajama validation set for next-token prediction as well as on downstream benchmarks that test abilities in common sense reasoning.

The third column in Table~\ref{tab:lmeval} displays the perplexities on the validation set for various approaches. Among the models, xLSTM[1:0] and xLSTM[7:1] achieve the lowest validation set perplexity across every model size. Performance on downstream tasks is presented in the remaining columns of Table~\ref{tab:lmeval}. Overall, xLSTM consistently outperforms competing methods for nearly all tasks and model sizes, with the exception that Mamba occasionally performs best on the ARC task. Additional information can be found in Appendix~\ref{sec:appExpLanguage}.

\paragraph{Performance on PALOMA Language Tasks.} Third, across each model size, we assess the next-token prediction abilities of xLSTM, RWKV-4, Llama, and Mamba using PALOMA language tasks~\citep{Magnusson:23arxivshort}. Evaluation is based on perplexity for next-token prediction across 571 diverse text domains, covering sources from nytimes.com to Reddit’s r/depression. 
Table~\ref{tab:ppleval} presents perplexity values, separating results for general language modeling tasks (first seven columns) from detailed, domain-specific benchmarks (final five columns). Notably, xLSTM[1:0] surpasses xLSTM[7:1] on these language-focused tasks.

xLSTM[1:0] demonstrates lower perplexity than Mamba in 568 of 571 (99.5\%) evaluated text domains, outperforms Llama in 486 of 571 (85.1\%) domains, and surpasses RWKV-4 in 570 of 571 (99.8\%) domains. Further information can be found in Appendix~\ref{sec:appExpLanguage}.

\paragraph{Scaling Laws.} Next, we investigate the scaling properties that adhere to a power-law relationship, which enable predictions of model performance as model size increases~\citep{Kaplan:20,Brown:20short}. Figure~\ref{fig:temp_sclaw300B} illustrates the respective scaling patterns. The analyzed models exhibit broadly similar power-law trends, yet are separated by distinct performance baselines. Among them, RWKV-4 demonstrates the weakest results, with Llama and Mamba subsequently outperforming it. xLSTM surpasses Mamba, achieving an advantage comparable to the gap between Mamba and Llama. These observed scaling patterns suggest that, as model sizes expand, xLSTM is poised to sustain a performance edge over both Transformer-based and State-Space models.

\textbf{Generation Times and Maximal Throughput.} We conclude by evaluating the generation latency in Figure~\ref{fig:inference_time_generation} and the peak throughput in Figure~\ref{fig:inference_time_decoding_throughput} (left) for our 1.3B-parameter xLSTM models. These results are compared against equivalently sized versions of Mamba, Llama, and RWKV from HuggingFace repositories, with the Llama model making use of a static key-value cache. During testing, compilation features such as full cache compilation for Transformers and \texttt{torch.compile} support for the Mamba model were unavailable. For all generation benchmarks, each model was run at batch size 1 with a pre-fill value of 16, a setting expected to be particularly advantageous for the Transformer architecture.

As depicted in Figure~\ref{fig:inference_time_generation}, both xLSTM and the other recurrent architectures (Mamba and RWKV-4) demonstrate linear growth in compute requirements, in contrast with Llama’s quadratic scaling. For the evaluation of decoding throughput, we varied batch sizes and prefill values specifically for the Llama model. The data in Figure~\ref{fig:inference_time_decoding_throughput} (right) illustrates that xLSTM’s constant memory footprint enables the use of much larger batch sizes than Llama, resulting in superior throughput performance.

\section{Limitations}

(i) Unlike mLSTM, the memory-mixing mechanism in sLSTM necessitates sequential operations, restricting the possibility of parallel execution and thus impeding the achievement of high-speed parallel implementations. Despite this, we have designed an efficient CUDA kernel for sLSTM, which, while currently operating at under twice the time required by our parallelized mLSTM variant, narrows the performance gap substantially. (ii) It should be noted that the current CUDA kernels for mLSTM lack optimization, making their execution approximately four times slower than highly optimized kernels such as FlashAttention or the scan operations implemented in Mamba. There is considerable room for improvement, and applying optimization techniques similar to those in FlashAttention may yield significant speed gains. (iii) The matrix memory component in mLSTM imposes elevated computational requirements owing to the necessity of handling $d \times d$ matrices. Nonetheless, both the memory update and retrieval processes are free of additional parameters and lend themselves well to standard matrix-based parallelization; as a result, the increased computational complexity has only a limited impact on actual runtime. (iv) Special attention should be given to the initialization of the forget gates, as this can critically affect performance. (v) Because the matrix memory does not scale with sequence length, increasing the length of the sequence has the potential to saturate the memory capacity when processing very long contexts. However, we have not observed this to be a constraint for context windows up to 16k tokens; refer to Section~\ref{sec:spaj300B_ctx_extrapolation} for further details. (vi) Owing to the significant computational expense involved in large-scale language model experimentation, neither the xLSTM architecture nor its hyperparameters have undergone comprehensive optimization to date, especially at larger model scales. We expect that rigorous architecture and hyperparameter tuning is necessary for xLSTM models to attain optimal performance.

\section{Conclusion}
Our investigation has partially addressed the question: To what extent can LSTM models achieve competitive results in language modeling when expanded to the billion-parameter scale? Our current findings suggest that, at a minimum, such scaling enables LSTMs to reach performance comparable to that of prevailing models, including Transformers and State Space Models. By introducing exponential gating, memory mixing, and an innovative memory structure, we have developed the xLSTM, an enhanced variant of the LSTM. Empirical results demonstrate that xLSTM achieves strong performance on language modeling benchmarks relative to leading architectures like Transformers and State Space Models. Analysis of scaling laws further suggests that as xLSTM models increase in size, they could emerge as formidable alternatives to contemporary Large Language Models based on Transformer architectures. Beyond language modeling, xLSTM may offer significant advancements in related domains such as Reinforcement Learning, Time Series Prediction, and modeling of physical processes.

\end{document}